% =========================================================================
\section{Circuito 2: amplificador no inversor}\label{sec:no-inversor}
% =========================================================================
\begin{prelab}[Amplificador no inversor]
Usando un modelo de amplificador operacional ideal, encuéntrese la ganancia en tensión
$A_V=\Vout/\Vin$ en función de las resistencias del circuito de la \cref{fig:no-inversor}. Si imponemos
$R_1=1\,\mathrm{k}\Omega$, determínese el valor de $R_2$ para conseguir $A_V=+11$ a $1\,\mathrm{kHz}$.
\end{prelab}

% Figura: amplificador no inversor (circuito 2) con tensiones y corrientes del análisis.
\begin{figure}[ht]
\centering
\begin{circuitikz}[circuito]
  \draw (0,0) node[op amp, ao, noinv input up] (A) {};
  % nudo inversor (abajo) y red de realimentación
  \draw (A.out) -- ++(0.7,0) coordinate (O);
  \draw (A.-) -- ++(-0.6,0) coordinate (N) -- ++(0,-1.0) coordinate (M);
  \draw (M) -- ++(-1.0,0) coordinate (G1) to[R, l_=$R_1$, i>^=$i_1$] ++(0,-1.7) node[ground] {};
  \draw (M) to[R, l_=$R_2$, i<_=$i_2$] (M -| O) -- (O);
  % generador en la entrada no inversora
  \draw (A.+) to[short, -o] ++(-2.0,0) coordinate (Ein) node[above, tens] {$\vin$};
  \draw (Ein) -- ++(-0.8,0) to[sV, l_=$v_g$] ++(0,-2.0) node[ground] {};
  % carga y salida
  \draw (O) to[short, *-o] ++(1.0,0) node[right, tens] {$\vout$};
  \draw (O) ++(0.5,0) to[R, l=$R_L$] ++(0,-2.0) node[ground] {};
  \fill (M) circle (1.6pt);
  \node[tens, anchor=east, font=\scriptsize] at ($(N)!0.5!(M)+(-0.08,0)$) {$\vm=\vin$};
\end{circuitikz}
\caption{Amplificador no inversor (figura~3 del guion), con $R_1=1\,\mathrm{k}\Omega$,
$R_L=10\,\mathrm{k}\Omega$ y la alimentación omitida. La entrada $+$ no absorbe corriente; el nudo
inversor sigue a la entrada ($\vm=\vin$) y $R_1$--$R_2$ forman un divisor de $\vout$.}
\label{fig:no-inversor}
\end{figure}


\begin{resolucion}
AO ideal con realimentación negativa a través de $R_2$; sentidos de corriente de la \cref{fig:no-inversor}.

\textbf{1. Cortocircuito virtual.} La señal entra por $+$ sin caída de tensión ($i_+=0$, R1); por (R2),
\begin{equation}\label{eq:n-vm}
  \vm=\vp=\vin .
\end{equation}

\textbf{2. Ley de Ohm en $R_1$}, entre $\vm$ y tierra:
\begin{equation}\label{eq:n-i1}
  i_1=\frac{\vm-0}{R_1}=\frac{\vin}{R_1} .
\end{equation}

\textbf{3. Kirchhoff en el nudo inversor}, con $i_-=0$ por (R1):
\begin{equation}\label{eq:n-i2}
  i_2=i_1+i_-\quad\Longrightarrow\quad i_2=i_1=\frac{\vin}{R_1} .
\end{equation}

\textbf{4. Ley de Ohm en $R_2$}, recorrida por $i_2$ de la salida al nudo inversor:
\begin{equation}\label{eq:n-vout}
  \vout-\vm=R_2\,i_2\quad\Longrightarrow\quad \vout=\vin+\frac{R_2}{R_1}\,\vin=\Bigl(1+\frac{R_2}{R_1}\Bigr)\vin .
\end{equation}

\textbf{Resultado.}
\begin{equation}\label{eq:n-Av}
  \resultado{A_V=\frac{\vout}{\vin}=1+\frac{R_2}{R_1}}
\end{equation}

\textbf{Diseño.} Con $A_V=11$ y $R_1=1\,\mathrm{k}\Omega$ en \eqref{eq:n-Av}:
\begin{equation}\label{eq:n-R2}
  R_2=(A_V-1)\,R_1=10\cdot1\,\mathrm{k}\Omega\quad\Longrightarrow\quad\resultado{R_2=10\,\mathrm{k}\Omega}
\end{equation}

$A_V>0$: salida en fase con la entrada; además $A_V\ge1$ siempre. Como en el inversor, la ganancia no
depende de la frecuencia, ni de $A$ ni de $R_L$.
\end{resolucion}
